\documentclass[11pt]{article}
\usepackage[a4paper,margin=1in]{geometry}
\usepackage{amsmath,amssymb,mathtools}
\usepackage{enumitem}
\usepackage{unicode-math}
\setlist{leftmargin=*}
\allowdisplaybreaks
\emergencystretch=3em
\title{KKT Endpoint-Bessel Reduction Criterion}
\author{}
\date{}
\begin{document}
\maketitle

Verdict: \textbf{A2-2 is useful; A3-2 is mostly not usable as a proof, but it points toward the right obstacle.} The best next step is not A3's shifted-\(Q\) Sonin functional. The better bridge is:

\begin{equation*}
\boxed{
\begin{gathered}
\text{direct }x\text{-Sonin localization}\\
+\ \text{regularized }W\text{-identity}\\
+\ \text{endpoint Bessel comparison with an adjustable frequency}
\end{gathered}
}
\end{equation*}


This stays analytic and does not use integer-parameter representation theory. The target remains the real-parameter extension of the KKT refined bound; the review notes that the refined estimate is known from the integer-parameter setting but the real-parameter version is the conjectural extension.  The same review identifies the Sonin--Polya/Sturm--Picone endpoint-Bessel route as the relevant analytic route.



\section*{1. What is actually useful in A2-2.tex}

A2-2's main useful observation is the regularized identity

\begin{equation*}
W(x):=F(x)g(x)^2+A(x)^2g'(x)^2,
\qquad A(x)=1-x^2,
\end{equation*}


with

\begin{equation*}
\boxed{
W'(x)=F'(x)g(x)^2.
}
\tag{1}
\end{equation*}


This is correct.

It follows directly from

\begin{equation*}
(A g')'+\frac{F}{A}g=0.
\end{equation*}


Indeed,

\begin{equation*}
\begin{aligned}
W'
&=F'g^2+2Fgg'+2AA'g'^2+2A^2g'g'' \\
&=F'g^2+2Fgg'+2AA'g'^2
+2g'(-AA'g'-Fg)\\
&=F'g^2.
\end{aligned}
\end{equation*}


This identity is better than the raw Sonin quotient near turning points because

\begin{equation*}
S_g(x)=g(x)^2+\frac{A(x)^2g'(x)^2}{F(x)}
\end{equation*}


blows up when \(F=0\), while

\begin{equation*}
\boxed{
W=F S_g
}
\end{equation*}


remains finite.

At the right endpoint,

\begin{equation*}
W(1)=0.
\end{equation*}


If \(\alpha>0\), this is because \(g(1)=0\) and \(A(1)=0\). If \(\alpha=0\), then \(F(1)=0\) and again \(W(1)=0\).

Therefore, if \(x_R\) is the first right endpoint extremum of \(g\), so that \(g'(x_R)=0\), then

\begin{equation*}
W(x_R)=F(x_R)g(x_R)^2,
\end{equation*}


and integrating (1) backward from (1) gives the exact tail identity

\begin{equation*}
\boxed{
g(x_R)^2
=
\frac{1}{F(x_R)}
\int_{x_R}^{1}(-F'(t))g(t)^2\,dt.
}
\tag{2}
\end{equation*}


This is the correct version of A2's ``boundary Volterra'' idea.

However, A2-2 stops at the right place: (2) alone does \textbf{not} close the proof. If one bounds \(g(t)^2\le g(x_R)^2\), then (2) only gives

\begin{equation*}
F(x_R)
\le
\int_{x_R}^{1}(-F'(t))\,dt
=
F(x_R)-F(1),
\end{equation*}


which is trivial because \(F(1)=-\alpha^2\le0\). The identity is useful as a regular bridge, not as a standalone bound.



\section*{2. What is wrong in A3-2.tex}

A3-2 has several mathematical defects.

First, it calls \(F\) convex. In this normalization,

\begin{equation*}
F(x)=\kappa(1-x^2)-\frac14(d-sx)^2
\end{equation*}


is a \textbf{concave} quadratic.

Second, A3 again confuses the Liouville variable with the original KKT function. The angular variable

\begin{equation*}
v(\theta)=\sqrt{\sin\theta},g(\cos\theta)
\end{equation*}


satisfies

\begin{equation*}
v''+Qv=0,
\end{equation*}


but the conjecture is about

\begin{equation*}
g(\cos\theta)=\frac{v(\theta)}{\sqrt{\sin\theta}}.
\end{equation*}


Uniformly controlling \(v\) does not uniformly control \(g\), especially at \(\theta=0,\pi\).

Third, A3's proposed shifted Sonin functional is not sign-controlled. If

\begin{equation*}
S_\lambda(\theta)=y^2+\frac{y'^2}{Q+\lambda},
\qquad y''+Qy=0,
\end{equation*}


then

\begin{equation*}
\boxed{
S_\lambda'
=
\frac{2\lambda yy'}{Q+\lambda}
-
\frac{Q'}{(Q+\lambda)^2}y'^2.
}
\tag{3}
\end{equation*}


The extra term

\begin{equation*}
\frac{2\lambda yy'}{Q+\lambda}
\end{equation*}


has no fixed sign. So the shifted-\(Q\) bridge is not a valid Sonin monotonicity argument.

Fourth, the statement that the (x)-functional is ``regular across turning points'' is only true after multiplying by \(F\). The quotient \(S_g\) still has \(1/F\). The regular object is

\begin{equation*}
W=F S_g,
\end{equation*}


not \(S_g\) itself.

So: \textbf{discard A3's proof mechanism; keep only its diagnosis that the turning-point singularity must be bridged.}



\section*{3. Correct direct equation and endpoint expansion}

Let

\begin{equation*}
u(x)=g_n^{(\alpha,\beta)}(x),
\qquad
s=\alpha+\beta,
\qquad
d=\beta-\alpha,
\end{equation*}


\begin{equation*}
A(x)=1-x^2,
\qquad
N=n+\frac{\alpha+\beta+1}{2},
\end{equation*}


\begin{equation*}
\kappa=n(n+s+1)+\frac{s}{2}.
\end{equation*}


Then

\begin{equation*}
\boxed{
(Au')'+G(x)u=0,
}
\tag{4}
\end{equation*}


where

\begin{equation*}
\boxed{
G(x)=
N^2-\frac14
-\frac{\alpha^2}{2(1-x)}
-\frac{\beta^2}{2(1+x)}.
}
\tag{5}
\end{equation*}


Equivalently,

\begin{equation*}
G(x)=\frac{F(x)}{A(x)},
\end{equation*}


with

\begin{equation*}
\boxed{
F(x)
=
\kappa(1-x^2)
-\frac14\bigl(\beta-\alpha-(\alpha+\beta)x\bigr)^2.
}
\tag{6}
\end{equation*}


Near the right endpoint, put

\begin{equation*}
x=1-\delta.
\end{equation*}


Then the exact quadratic expansion is

\begin{equation*}
\boxed{
F(1-\delta)
=
-\alpha^2
+
L_R\delta
-
\left(N^2-\frac14\right)\delta^2,
}
\tag{7}
\end{equation*}


where

\begin{equation*}
\boxed{
L_R=2\kappa+\alpha(\alpha+\beta)
=
2N^2-\frac12+\frac{\alpha^2-\beta^2}{2}.
}
\tag{8}
\end{equation*}


Similarly, near the left endpoint,

\begin{equation*}
\boxed{
F(-1+\delta)
=
-\beta^2
+
L_L\delta
-
\left(N^2-\frac14\right)\delta^2,
}
\tag{9}
\end{equation*}


where

\begin{equation*}
\boxed{
L_L=2\kappa+\beta(\alpha+\beta)
=
2N^2-\frac12+\frac{\beta^2-\alpha^2}{2}.
}
\tag{10}
\end{equation*}


This explains the apparent \(0/0\) problem in A2. When \(\alpha\downarrow0\), the right turning point satisfies approximately

\begin{equation*}
\delta_R\sim \frac{\alpha^2}{L_R},
\end{equation*}


so both numerator and denominator in the Volterra identity collapse at the Bessel scale. The singularity is not fatal; it says that the correct endpoint variable is Bessel-type.



\section*{4. Direct Sonin localization remains the global skeleton}

On \(F>0\),

\begin{equation*}
\boxed{
S_g(x)=u(x)^2+\frac{A(x)^2u'(x)^2}{F(x)}.
}
\tag{11}
\end{equation*}


The exact identity is

\begin{equation*}
\boxed{
S_g'(x)
=
-\frac{A(x)^2F'(x)}{F(x)^2}u'(x)^2.
}
\tag{12}
\end{equation*}


Because \(F\) is concave, with vertex \(x_0\),

\begin{equation*}
S_g \text{ decreases on }(x_-,x_0),
\end{equation*}


and

\begin{equation*}
S_g \text{ increases on }(x_0,x_+).
\end{equation*}


At a critical point of \(u\),

\begin{equation*}
u'=0,
\qquad
S_g=u^2.
\end{equation*}


Therefore the global maximum of \(|u|\) is reduced to the first endpoint lobe near \(x=1\) and the first endpoint lobe near \(x=-1\).

So the remaining problem is exactly:

\begin{equation*}
\boxed{
\text{prove the sharp KKT bound on each first endpoint lobe.}
}
\end{equation*}




\section*{5. Corrected endpoint-Bessel comparison with adjustable frequency}

Set

\begin{equation*}
x=\cos\theta,
\qquad
0<\theta<\pi,
\end{equation*}


and

\begin{equation*}
v(\theta)=\sqrt{\sin\theta},u(\cos\theta).
\end{equation*}


Then

\begin{equation*}
\boxed{
v''(\theta)+Q(\theta)v(\theta)=0,
}
\tag{13}
\end{equation*}


where

\begin{equation*}
\boxed{
Q(\theta)
=
N^2
+
\frac{1/4-\alpha^2}{4\sin^2(\theta/2)}
+
\frac{1/4-\beta^2}{4\cos^2(\theta/2)}.
}
\tag{14}
\end{equation*}


The earlier small-parameter proof used the crude comparison

\begin{equation*}
Q(\theta)\ge
N^2+\frac{1/4-\alpha^2}{\theta^2},
\end{equation*}


which works only when

\begin{equation*}
0\le\alpha,\beta\le\frac12.
\end{equation*}


To go further, replace \(N\) by an \textbf{adjustable lower comparison frequency} \(M_R\).

Fix \(0<\Theta<\pi\). Define

\begin{equation*}
c_\alpha=\frac14-\alpha^2,
\qquad
c_\beta=\frac14-\beta^2,
\end{equation*}


\begin{equation*}
B(\theta)=\frac{1}{4\sin^2(\theta/2)}-\frac{1}{\theta^2},
\end{equation*}


\begin{equation*}
C(\theta)=\frac{1}{4\cos^2(\theta/2)}.
\end{equation*}


Then

\begin{equation*}
Q(\theta)
-
\frac{c_\alpha}{\theta^2}
=
N^2+c_\alpha B(\theta)+c_\beta C(\theta).
\end{equation*}


On \(0<\theta\le\Theta\), (B) and (C) are increasing. Hence

\begin{equation*}
Q(\theta)
\ge
M_R(\Theta)^2+\frac{1/4-\alpha^2}{\theta^2},
\tag{15}
\end{equation*}


where

\begin{equation*}
\boxed{
M_R(\Theta)^2
=
N^2+m_\alpha(\Theta)+m_\beta(\Theta),
}
\tag{16}
\end{equation*}


with

\begin{equation*}
m_\alpha(\Theta)=
\begin{cases}
\dfrac{1/4-\alpha^2}{12},
&0\le\alpha\le\frac12,\\[1.2ex]
\left(\dfrac14-\alpha^2\right)B(\Theta),
&\alpha>\frac12,
\end{cases}
\tag{17}
\end{equation*}


and

\begin{equation*}
m_\beta(\Theta)=
\begin{cases}
\dfrac{1/4-\beta^2}{4},
&0\le\beta\le\frac12,\\[1.2ex]
\left(\dfrac14-\beta^2\right)C(\Theta),
&\beta>\frac12.
\end{cases}
\tag{18}
\end{equation*}


This is the corrected version of the endpoint comparison. It permits \(\beta>1/2\), and even \(\alpha>1/2\), provided \(M_R(\Theta)^2\) remains positive and large enough.



\section*{6. Endpoint-lobe criterion}

Let \(j_{\alpha,1}\) be the first positive zero of \(J_\alpha\). Suppose there is some \(0<\Theta<\pi\) such that

\begin{equation*}
\boxed{
M_R(\Theta)>0,
\qquad
\frac{j_{\alpha,1}}{M_R(\Theta)}\le \Theta.
}
\tag{19}
\end{equation*}


Then on the first right endpoint lobe, Sturm comparison gives

\begin{equation*}
0\le v(\theta)
\le
K_R\sqrt{\theta},J_\alpha(M_R\theta),
\tag{20}
\end{equation*}


where matching the regular endpoint asymptotic gives

\begin{equation*}
\boxed{
K_R
=
C_{n,\alpha,\beta}
\frac{\Gamma(n+\alpha+1)}
{\Gamma(n+1)M_R^\alpha}.
}
\tag{21}
\end{equation*}


Since

\begin{equation*}
u(\cos\theta)=\frac{v(\theta)}{\sqrt{\sin\theta}},
\end{equation*}


we get

\begin{equation*}
\boxed{
|u(\cos\theta)|
\le
K_R
\left(\frac{\theta}{\sin\theta}\right)^{1/2}
J_\alpha(M_R\theta)
}
\tag{22}
\end{equation*}


through the first right endpoint lobe.

Now define

\begin{equation*}
\rho(\Theta)
=
\sup_{0<t\le\Theta}
\frac{1}{2t^2}\log\frac{t}{\sin t}.
\end{equation*}


Then

\begin{equation*}
\left(\frac{\theta}{\sin\theta}\right)^{1/2}
\le
e^{\rho(\Theta)\theta^2}.
\end{equation*}


Also, for \(0\le z\le j_{\alpha,1}\),

\begin{equation*}
J_\alpha(z)
\le
\frac{(z/2)^\alpha}{\Gamma(\alpha+1)}
\exp\left(-\frac{z^2}{4(\alpha+1)}\right).
\end{equation*}


Therefore, if

\begin{equation*}
\boxed{
c_R(\Theta)
:=
\frac{M_R(\Theta)^2}{4(\alpha+1)}
-
\rho(\Theta)
> 0,
}
\tag{23}
\end{equation*}


then

\begin{equation*}
|u(\cos\theta)|
\le
D_R
\theta^\alpha e^{-c_R\theta^2},
\tag{24}
\end{equation*}


where

\begin{equation*}
\boxed{
D_R
=
\frac{
C_{n,\alpha,\beta}\Gamma(n+\alpha+1)
}
{
\Gamma(n+1)2^\alpha\Gamma(\alpha+1)
}.
}
\tag{25}
\end{equation*}


Optimizing the scalar function \(\theta^\alpha e^{-c_R\theta^2}\) gives

\begin{equation*}
\max_{\theta\ge0}\theta^\alpha e^{-c_R\theta^2}
=
\left(\frac{\alpha}{2e c_R}\right)^{\alpha/2},
\end{equation*}


with value (1) when \(\alpha=0\).

Thus the first right endpoint lobe is bounded by

\begin{equation*}
\boxed{
E_R(\Theta)
=
D_R
\left(\frac{\alpha}{2e c_R(\Theta)}\right)^{\alpha/2}.
}
\tag{26}
\end{equation*}


The corresponding left endpoint criterion is obtained by swapping

\begin{equation*}
\alpha\leftrightarrow\beta.
\end{equation*}


Let

\begin{equation*}
E_L(\Theta)
=
E_R(\Theta)
\quad\text{with } \alpha,\beta \text{ swapped.}
\end{equation*}




\section*{7. New reduction theorem}

Define the KKT target constant

\begin{equation*}
\boxed{
T_{n,\alpha,\beta}
=
\left(
\frac{(n+1)(n+\alpha+\beta+1)}
{(n+\alpha+1)(n+\beta+1)}
\right)^{1/4}.
}
\tag{27}
\end{equation*}


Then the following is a clean reduction.

\subsection*{Proposition}

Assume that for the right endpoint there exists \(\Theta_R\in(0,\pi)\) satisfying

\begin{equation*}
M_R(\Theta_R)>0,
\qquad
\frac{j_{\alpha,1}}{M_R(\Theta_R)}\le\Theta_R,
\qquad
c_R(\Theta_R)>0,
\end{equation*}


and

\begin{equation*}
\boxed{
E_R(\Theta_R)\le T_{n,\alpha,\beta}.
}
\tag{28}
\end{equation*}


Assume also that for the left endpoint, after swapping \(\alpha,\beta\), there exists \(\Theta_L\in(0,\pi)\) satisfying the analogous conditions and

\begin{equation*}
\boxed{
E_L(\Theta_L)\le T_{n,\alpha,\beta}.
}
\tag{29}
\end{equation*}


Then

\begin{equation*}
\boxed{
|g_n^{(\alpha,\beta)}(x)|
\le
T_{n,\alpha,\beta}
\qquad (-1\le x\le1).
}
\tag{30}
\end{equation*}


\subsection*{Proof}

The endpoint comparison proves the target bound on the first right and first left endpoint lobes. The direct (x)-Sonin functional proves that all interior local extrema are no larger than the corresponding first endpoint-lobe extrema. Endpoint values are harmless: if \(\alpha=0\), then \(g(1)=1\) and \(T_{n,0,\beta}=1\); if \(\beta=0\), then \(|g(-1)|=1\) and \(T_{n,\alpha,0}=1\). Hence the global bound follows.

This is a genuine advance over A2/A3: the turning-point singularity is no longer attacked by a false shifted Sonin functional. It is bypassed by a localized Bessel comparison whose constants are explicit.



\section*{8. Why this improves the previous small-parameter proof}

For

\begin{equation*}
0\le\alpha,\beta\le\frac12,
\end{equation*}


one can simply take

\begin{equation*}
M_R=N
\end{equation*}


because both inverse-square endpoint corrections in \(Q\) are nonnegative. With \(\Theta=2\pi/3\), one has

\begin{equation*}
\rho(\Theta)\le \frac18,
\end{equation*}


and the previous small-square proof follows.

The new criterion is stronger because when, say, \(\beta>1/2\), the negative contribution

\begin{equation*}
\frac{1/4-\beta^2}{4\cos^2(\theta/2)}
\end{equation*}


is not treated as fatal. It is absorbed into the lowered frequency \(M_R(\Theta)\). This is exactly the missing mechanism outside the square

\begin{equation*}
0\le\alpha,\beta\le\frac12.
\end{equation*}




\section*{9. Remaining hard work}

The problem is now reduced to proving explicit inequalities of the form

\begin{equation*}
E_R(\Theta_R)\le T_{n,\alpha,\beta},
\qquad
E_L(\Theta_L)\le T_{n,\alpha,\beta}.
\end{equation*}


There are three plausible ways to finish:

1. \textbf{Analytic optimization in \(\Theta\).}
   Choose \(\Theta\) as a simple multiple of \(j_{\alpha,1}/M_R\). This should sharpen \(M_R\) while keeping \(\rho(\Theta)\) small.

2. \textbf{Bessel-zero inequalities.}
   Use standard bounds for \(j_{\nu,1}\), such as
\begin{equation*}
j_{\nu,1}<\nu+c\nu^{1/3}+C,
\end{equation*}

   to make condition (19) algebraic.

3. \textbf{Hybrid interval proof.}
   The large-balanced regime \(\alpha,\beta\ge2n\) is already handled by the weighted-energy argument. The small square is already handled by the \(M=N\) Bessel comparison. The remaining middle region can be attacked by verifying (28)--(29) with rigorous interval arithmetic, because all quantities are explicit gamma/Bessel/trigonometric expressions.



\section*{Bottom line}

A2-2 gives a useful exact identity:

\begin{equation*}
\boxed{
W=F g^2+A^2g'^2,\qquad W'=F'g^2.
}
\end{equation*}


A3-2's proposed shifted Sonin bridge is not valid because the derivative acquires an uncontrolled term.

The best forward route is the endpoint criterion above. It converts the remaining turning-point problem into explicit inequalities involving

\begin{equation*}
M_R(\Theta),\quad M_L(\Theta),\quad j_{\alpha,1},\quad j_{\beta,1},
\quad
\Gamma\text{-ratios}.
\end{equation*}


This is now a concrete analytic/finite-verification path rather than a qualitative strategy.
\end{document}
